\chapter*{Resumo}
\label{ch:resumo}

A crescente geração de resíduos alimentares e o excedente de glicerol bruto proveniente da produção de biodiesel representam duas produções significativas de resíduos orgânicos que requerem soluções sustentáveis. Este estudo avaliou a viabilidade técnica de um sistema de codigestão anaeróbia de uma única etapa, tratando resíduos alimentares e glicerol como co-substratos, com o objetivo de verificar a produção de biogás e, simultaneamente, proporcionar uma via de valorização para o glicerol bruto. A investigação foi conduzida num reator contínuo de tanque agitado de 1 L, operado em condições mesofílicas (35 °C), com um tempo de retenção hidráulica de 30 dias. Os resíduos alimentares foram recolhidos numa cantina universitária e o glicerol bruto foi obtido a partir de resíduos da produção de biodiesel; lamas anaeróbias de uma estação de tratamento de águas residuais municipais serviram como inóculo. O reator foi operado em fases sequenciais: um período inicial de adaptação, codigestão de resíduos alimentares e glicerol, alimentação exclusiva com glicerol e um período pós-alimentação. A produção de biogás foi monitorizada continuamente através de um \textit{automatic methane potential test system}, enquanto a estabilidade do processo foi avaliada por medições de pH, alcalinidade total, ácidos graxos voláteis e respectiva razão. A caracterização dos substratos incluiu análises de sólidos totais, sólidos voláteis, densidade, carbono total e nitrogênio total. Os resultados demonstraram que a codigestão de resíduos alimentares e glicerol originou as taxas de produção de metano e os volumes acumulados mais elevados, com picos imediatamente após as alimentações, indicando uma rápida conversão do substrato. Em contraste, a alimentação exclusiva com glicerol resultou numa produção reduzida de biogás, apesar da atividade metanogénica continuada. A alcalinidade total manteve-se elevada após a suplementação com bicarbonato, mantendo o pH em torno de 7 durante a maior parte da experiência; as concentrações de ácidos graxos voláteis flutuaram entre 3000 e 5000 mg CH\textsubscript{3}COOH/L, com razões AGV/alcalinidade de 0,3 a 1,2. Estes resultados confirmam a hipótese de que o glicerol pode ser útil para a produção de biogás quando codigerido com resíduos alimentares ou isoladamente. O estudo demonstra a viabilidade técnica da codigestão anaeróbia como estratégia de valorização do glicerol bruto no âmbito de uma bioeconomia circular, embora seja necessária investigação adicional para otimizar as cargas de glicerol e desenvolver estratégias de alimentação robustas para aplicação industrial.
